\newpage % The Meta-Review should at least start on a new column

% Use \appendices and not \appendix due to IEEEtran.cls quirks
%\appendices % if not used earlier

% \section{Meta-Review}

% The following meta-review was prepared by the program committee for the 2026
% IEEE Symposium on Security and Privacy (S\&P) as part of the review process as
% detailed in the call for papers.

% \subsection{Summary}
% This paper investigates poisoning attacks on GraphRAG, a retrieval-augmented generation method that builds knowledge graphs from text. The authors introduce GRAGPoison, an attack that exploits shared relations in the knowledge graph to poison multiple queries simultaneously. GRAGPoison combines relation injection, relation enhancement, and narrative generation to craft poisoning texts. The evaluation spans several datasets (geographic, medical, cyber-security, and MuSiQue) and includes ablation studies and preliminary defense exploration.

% \subsection{Scientific Contributions}
% \begin{itemize}
% \item New attack design
% \item Empirical evidence of vulnerability
% \item New research direction
% \end{itemize}

% \subsection{Reasons for Acceptance}
% \begin{enumerate}
% \item The problem is timely and important: GraphRAG is growing in adoption, yet its robustness has not been systematically evaluated.
% \item GRAGPoison is a novel and conceptually simple attack that demonstrates clear vulnerabilities in GraphRAG.
% \item The empirical evaluation shows strong attack performance and includes ablations to analyze design choices.
% \item The work provides resources (datasets and methodology) that can be valuable for the community in reproducing and extending the findings.
% \end{enumerate}